% test-autoref.tex
% 自動基準（既定で現在フォントの x ハイトを基準）のストレステスト。
% コンパイル: latexmk -r .latexmkrc test-autoref.tex   (LuaLaTeX 専用)
\documentclass{jlreq}
\usepackage[tsumesuji=1]{KKsymbols}
\usepackage{luatexja-otf}
\IfFileExists{KKran.sty}{\usepackage{KKran}}{\newcommand{\Rrnum}[1]{\textit{\romannumeral #1}}}

% ===== フォント一式（mc/gt/mg を別フォントに割り当てて font 感応をテスト） =====
\makeatletter
\RequirePackage[no-math]{fontspec}
\RequirePackage[no-math,match,scale=1]{luatexja-fontspec}
\RequirePackage[hiragino-pro,deluxe,expert]{luatexja-preset}
\setmainfont{HiraMinPro-W3}[BoldFont=HiraMinPro-W6]
\setmainjfont{HiraMinPro-W3}[BoldFont=HiraMinPro-W6]
\newfontfamily{\sfhira@pre}{HiraKakuPro-W3}[BoldFont=HiraKakuPro-W6]
\newjfontfamily{\sfhiraj@pre}{HiraKakuPro-W3}[BoldFont=HiraKakuPro-W6]
\newfontfamily{\mchira@pre}{HiraMinPro-W3}[BoldFont=HiraMinPro-W6]
\newjfontfamily{\mchiraj@pre}{HiraMinPro-W3}[BoldFont=HiraMinPro-W6]
\newfontfamily{\gthira@pre}{HiraKakuPro-W3}[BoldFont=HiraKakuPro-W6,FontFace={eb}{\shapedefault}{HiraKakuStd-W8}]
\newjfontfamily{\gthiraj@pre}{HiraKakuPro-W3}[BoldFont=HiraKakuPro-W6,FontFace={eb}{\shapedefault}{HiraKakuStd-W8}]
\newfontfamily{\mghira@pre}{HiraMaruPro-W4}
\newjfontfamily{\mghiraj@pre}{HiraMaruPro-W4}
\renewcommand{\sffamily}{\sfhira@pre\sfhiraj@pre}
\renewcommand{\mcfamily}{\mchira@pre\mchiraj@pre}
\renewcommand{\gtfamily}{\gthira@pre\gthiraj@pre}
\renewcommand{\mgfamily}{\mghira@pre\mghiraj@pre}
\makeatother

\fboxsep=0pt\fboxrule=.1pt
\newcommand{\romanrow}{%
  \fbox{\maru{\Rrnum{1}}}\fbox{\maru{\Rrnum{2}}}\fbox{\maru{\Rrnum{3}}}%
  \fbox{\maru{\Rrnum{4}}}\fbox{\maru{\Rrnum{6}}}\fbox{\maru{\Rrnum{7}}}\fbox{\maru{\Rrnum{8}}}}

\begin{document}

\section*{1. 自動基準: 丸囲みローマ数字（何も書かない）}
\noindent\textbf{(a) \texttt{\textbackslash kksrefoff}（旧来＝引数自身の寸法。バラつく）:}\par
{\kksrefoff\romanrow}
\par\medskip\noindent\textbf{(b) 手書き \texttt{\textbackslash vphantom\{\textbackslash Rrnum\{6\}\}}（従来）:}\par
\fbox{\maru{\vphantom{\Rrnum{6}}\Rrnum{1}}}\fbox{\maru{\vphantom{\Rrnum{6}}\Rrnum{2}}}%
\fbox{\maru{\vphantom{\Rrnum{6}}\Rrnum{3}}}\fbox{\maru{\vphantom{\Rrnum{6}}\Rrnum{4}}}%
\fbox{\maru{\vphantom{\Rrnum{6}}\Rrnum{6}}}\fbox{\maru{\vphantom{\Rrnum{6}}\Rrnum{7}}}%
\fbox{\maru{\vphantom{\Rrnum{6}}\Rrnum{8}}}
\par\medskip\noindent\textbf{(c) 既定＝自動（\texttt{\textbackslash maru\{\textbackslash Rrnum\{n\}\}} のみ）:}\par
\romanrow

\section*{2. フォント感応（mc/gt/mg で字形・x ハイトが変わる）}
\noindent\textbf{明朝(mc):} {\mcfamily \fbox{\maru{\Rrnum{3}}}\fbox{\maru{8}}\fbox{\maru{A}}\fbox{\maru{あ}}}\par
\noindent\textbf{ゴシック(gt):} {\gtfamily \fbox{\maru{\Rrnum{3}}}\fbox{\maru{8}}\fbox{\maru{A}}\fbox{\maru{あ}}}\par
\noindent\textbf{丸ゴ(mg):} {\mgfamily \fbox{\maru{\Rrnum{3}}}\fbox{\maru{8}}\fbox{\maru{A}}\fbox{\maru{あ}}}

\section*{3. \texttt{\textbackslash Rrnum} と他の内容が同居するケース}
\noindent 期待: 背の高い字が混ざれば自動基準は no-op（その字の寸法が支配）。\par
\fbox{\maru{A\Rrnum{3}}}\fbox{\maru{\Rrnum{3}A}}\fbox{\maru{\Rrnum{3}g}}%
\fbox{\maru{(\Rrnum{3})}}\fbox{\maru{\Rrnum{12}}}\fbox{\maru{第\Rrnum{3}}}%
\fbox{\maru{\Rrnum{3}.}}\fbox{\maru{\textbf{\Rrnum{3}}}}\fbox{\maru{\Rrnum{3}\Rrnum{8}}}

\section*{4. 純小文字（自動基準が効く）}
\noindent\fbox{\maru{a}}\fbox{\maru{c}}\fbox{\maru{e}}\fbox{\maru{m}}\fbox{\maru{x}}\fbox{\maru{n}}\quad
\textbf{[参考] 旧来:} {\kksrefoff\fbox{\maru{a}}\fbox{\maru{c}}\fbox{\maru{e}}\fbox{\maru{m}}\fbox{\maru{x}}\fbox{\maru{n}}}

\section*{5. 後方互換（背の高い内容は不変のはず）}
\noindent\maru{1}\maru{97}\maru{だ}\maru{ばばば}\maru{Qjg}\maru{A}\kuromaru{3}\nmaru{亀}\seihou{8}\kakko{12}

\end{document}
